\section{ALNS con Q-learning}
\label{sec:ql}

\cod{ALNS\_QLearning} sustituye las dos ruletas por un agente de aprendizaje por refuerzo
tabular que aplica Q-learning~\cite{watkins1992,sutton2018}. La elección de operadores se
modela como un proceso de decisión con los siguientes elementos.

\begin{itemize}
\item \textbf{Estados.} Dos estados, $z\in\{0,1\}$: $z=1$ si la iteración anterior produjo
una mejora y $z=0$ en caso contrario. El estado inicial es $z=0$.
\item \textbf{Acciones.} Cada acción es un \emph{par} $\theta=(j,k)$ de operador de
destrucción y de reparación, codificado como $\theta=5j+k$. Hay $6\times 5=30$ acciones.
\item \textbf{Tabla.} $Q\in\mathbb{R}^{2\times 30}$, inicializada en cero.
\item \textbf{Memoria de mejoras.} Un vector $M\in\mathbb{R}^{30}$, inicializado en cero, con
la mejora reciente de cada acción (\cod{recent\_improvement}); se usa para calcular el costo
de oportunidad (Sección~\ref{sec:recompensa}).
\end{itemize}

A diferencia del método clásico, la unidad de decisión es el par completo: el agente puede
aprender que una destrucción funciona con una reparación y no con otra. A cambio debe estimar
$60$ valores en lugar de $11$ pesos.

\begin{table}[H]
\centering
\small
\caption{Parámetros del Q-learning (\cod{SolverParams}, valores por defecto). Todos pueden
modificarse desde la línea de comandos con la sintaxis \cod{clave=valor}. La última fila es
una constante de \cod{alns\_qlearning.cpp}, no un parámetro.}
\label{tab:params-ql}
\begin{tabularx}{\textwidth}{@{}l c c L@{}}
\toprule
Parámetro & Símbolo & Valor & Función \\
\midrule
\cod{alpha} & $\alpha$ & $0.5$ & Tasa de aprendizaje \\
\cod{gamma} & $\gamma$ & $0.7$ & Factor de descuento \\
\cod{learning\_loop} & $T_{\mathrm{exp}}$ & $200$ & Iteraciones de exploración pura al inicio de la búsqueda y al inicio de la fase~2 \\
\cod{beta} & $\beta$ & $0.99$ & Decaimiento multiplicativo de $\varepsilon$ por iteración \\
\cod{epsilon\_min} & $\varepsilon_{\min}$ & $0.01$ & Exploración residual \\
\cod{eta} & $\eta$ & $0.8$ & Peso de la mejora global frente a la local en la recompensa \\
\cod{reward\_gain} & $\kappa$ & $10^4$ & Ganancia de la transformación logarítmica \\
\midrule
\cod{IMPROVEMENT\_DECAY} & $\delta$ & $0.9$ & Inercia de la memoria de mejoras $M$ \\
\bottomrule
\end{tabularx}
\end{table}

\subsection{Selección de la acción}

La política es $\varepsilon$-voraz con un calendario de exploración en tres tramos
(Algoritmo~\ref{alg:ql-sel}). Sea $t_0$ la iteración en la que comienza el calendario
($t_0=0$ al inicio de la búsqueda) y $t_{\mathrm{exp}}=t_0+T_{\mathrm{exp}}$ la última
iteración de exploración pura (\cod{explore\_until}):
\begin{equation}
\varepsilon_t=
\begin{cases}
1 & t\le t_{\mathrm{exp}},\\[2pt]
\max\bigl(\varepsilon_{\min},\ \beta^{\,t-t_{\mathrm{exp}}-1}\bigr) & t>t_{\mathrm{exp}}.
\end{cases}
\label{eq:eps}
\end{equation}

\begin{enumerate}
\item \emph{Calentamiento} (las primeras $200$ iteraciones): las acciones se eligen
uniformemente al azar. Cada par se prueba en promedio $200/30\approx 7$ veces y la tabla se
actualiza con normalidad, de modo que al terminar ya contiene una primera estimación de todos
los pares.
\item \emph{Transición} (las $459$ siguientes): $\varepsilon$ decae geométricamente desde $1$
y llega a $\varepsilon_{\min}$, pues $\lceil\ln 0.01/\ln 0.99\rceil=459$.
\item \emph{Explotación} (a partir de la iteración $660$ del calendario, aproximadamente):
$\varepsilon=0.01$ y el 99\,\% de las decisiones son voraces.
\end{enumerate}

El calendario se recorre \emph{dos veces}: una desde el inicio de la búsqueda y otra desde el
inicio de la fase~2 (Sección~\ref{sec:ql-fase2}). En la fase~2, que dura
$T_{\max}=25\,000$ iteraciones, más del 97\,\% transcurre en explotación. La fase~1 tiene
duración variable; si termina antes de unas $660$ iteraciones, su calendario queda truncado y
toda ella transcurre en calentamiento o transición.

En la rama voraz se elige la acción de mayor valor en el estado actual. Los empates (con
tolerancia $10^{-12}$) se resuelven al azar, lo que evita un sesgo sistemático hacia las
acciones de índice bajo mientras la tabla conserva valores iguales.

\begin{algorithm}[H]
\caption{Selección $\varepsilon$-voraz (\cod{ALNS\_QLearning::selectOperators})}
\label{alg:ql-sel}
\Entrada{iteración $t$; estado $z$; tabla $Q$; exploración actual $\varepsilon$; fin de la exploración pura $t_{\mathrm{exp}}$}
\Salida{par $(j,k)$; se memoriza la acción $\theta$}
$\varepsilon_t\gets 1$ si $t\le t_{\mathrm{exp}}$; si no, $\varepsilon$\;
\uSi{$\unif{0}{1}<\varepsilon_t$}{
  $\theta\sim\unifint{0}{29}$ \Com*[r]{exploración}
}
\SiNo{
  $Q^\star\gets\max_{\theta'}Q(z,\theta')$\;
  $\theta\gets$ elemento al azar de $\{\theta':Q(z,\theta')\ge Q^\star-10^{-12}\}$ \Com*[r]{explotación}
}
\Devolver $(\lfloor\theta/5\rfloor,\ \theta\bmod 5)$\;
\end{algorithm}

\subsection{Función de recompensa}
\label{sec:recompensa}

La recompensa se construye a partir de las dos ganancias relativas que calcula el bucle común
(Sección~\ref{sec:ganancia}), ambas medidas \emph{antes} de actualizar $x$ y $\xbest$:
\[
\gbest_t=G(\xbest,x'),\qquad \gcur_t=G(x,x'),
\]
con $G$ definida en~\eqref{eq:ganancia}: la mejora relativa en el primer criterio (clientes
sin asignar, vehículos o distancia) en que la candidata difiere de la referencia.

\paragraph{1. Transformación logarítmica.} Cada ganancia se transforma con
\begin{equation}
\varphi(g)=\ln(1+\kappa\,g)\ \text{ si } g>0,\qquad \varphi(g)=0\ \text{ si } g\le 0 .
\label{eq:phi}
\end{equation}
Su papel se entiende a partir de las escalas de $G$. Una reducción de distancia $\Delta D$
produce $g=\Delta D/D$, un valor pequeño (por ejemplo, $10^{-3}$ al ahorrar $5$ unidades
sobre $5000$); ahorrar un vehículo produce
$g=1/K$ y colocar $c$ clientes del banco produce $g=c/|\Uns|$, valores uno o varios órdenes de
magnitud mayores. Con $\kappa=10^{4}$ se obtiene
\[
\varphi=\ln\Bigl(1+10^{4}\,\frac{\Delta D}{D}\Bigr),\qquad
\varphi=\ln\Bigl(1+\frac{10^{4}}{K}\Bigr),\qquad
\varphi=\ln\Bigl(1+10^{4}\,\frac{c}{|\Uns|}\Bigr),
\]
respectivamente. La ganancia $\kappa$ fija la resolución de la señal: las mejoras relativas
menores que $1/\kappa=10^{-4}$ quedan en el tramo casi lineal de $\varphi$ y las mayores se
comprimen. El logaritmo impide así que los avances en los criterios prioritarios sepulten la
señal de las mejoras de distancia, sin dejar de ordenarlos: por ejemplo, con $K=20$ y
$D=5000$, ahorrar $10$ unidades de distancia da $\varphi\approx 3.0$, eliminar un vehículo da
$\varphi\approx 6.2$ y colocar uno de $8$ clientes pendientes da $\varphi\approx 7.1$. Las
ganancias por clientes colocados son las características de la fase~1; en la fase~2
predominan las de distancia.

\paragraph{2. Combinación global--local.} Las dos ganancias transformadas se mezclan con peso
$\eta=0.8$:
\begin{equation}
I_t=\eta\,\varphi(\gbest_t)+(1-\eta)\,\varphi(\gcur_t).
\label{eq:mejora}
\end{equation}
Como $\xbest$ nunca es peor que $x$, una ganancia global positiva implica una ganancia local
positiva; por tanto $I_t>0$ si y solo si la candidata mejora a la solución actual. Una mejora
solo local recibe el 20\,\% de su valor transformado; un nuevo mejor recibe además el 80\,\%
de $\varphi(\gbest_t)$. El agente es recompensado por cualquier progreso, pero sobre todo por
superar el mejor conocido.

\paragraph{3. Escalado temporal.} La recompensa se multiplica por la fracción de búsqueda
transcurrida, $\tau_t=t/T_{\mathrm{fin}}\in(0,1]$, donde $T_{\mathrm{fin}}$ es la última
iteración prevista en ese momento (Algoritmo~\ref{alg:alns}). Las mejoras son abundantes al
principio y escasas al final; el factor $\tau_t$ da más valor a las mejoras tardías, que son
las difíciles, y reduce el peso de las tempranas, que casi cualquier operador consigue. Como
$T_{\mathrm{fin}}=2\,T_{\max}$ durante la fase~1 y $T_{\mathrm{fin}}=t_1+T_{\max}$ durante la
fase~2 (con $t_1$ la iteración en que acabó la fase~1), $\tau_t$ no supera $0.5$ en la fase~1
y crece desde $t_1/(t_1+T_{\max})$ hasta $1$ en la fase~2.

\paragraph{4. Costo de oportunidad.} El agente mantiene, para cada acción $\theta$, una media
móvil exponencial $M_\theta$ de la mejora $I$ que ha producido en sus usos recientes, contando
como $0$ los usos sin mejora. Una iteración sin mejora no recibe recompensa nula sino
negativa, igual a la mayor mejora reciente \emph{de las demás acciones}:
\begin{equation}
r_t=
\begin{cases}
\tau_t\,I_t & \text{si } I_t>0,\\[2pt]
-\,\tau_t\,\Omega_t & \text{si } I_t=0,
\end{cases}
\qquad\text{con}\quad \Omega_t=\max_{\theta'\neq\theta}M_{\theta'} .
\label{eq:recompensa}
\end{equation}
Calculada la recompensa, se actualiza la memoria de la acción usada, haya o no mejora:
\begin{equation}
M_\theta\gets\delta\,M_\theta+(1-\delta)\,I_t,\qquad \delta=0.9 .
\label{eq:memoria}
\end{equation}
La interpretación es que una iteración desperdiciada cuesta lo que se habría podido ganar con
la mejor alternativa. $M_\theta$ estima la mejora \emph{media por uso} de la acción $\theta$
sobre una ventana efectiva de unos $1/(1-\delta)=10$ usos, de modo que combina con qué
frecuencia mejora y cuánto. De la definición se siguen tres propiedades:
\begin{itemize}
\item la penalización es \emph{adaptativa}: sube cuando alguna otra acción está produciendo
mejoras y se anula cuando ninguna lo hace, en cuyo caso fallar no cuesta nada;
\item al excluir a la propia acción del máximo, un par no se ve penalizado por sus propios
éxitos anteriores, sino solo por lo que ofrecen las alternativas;
\item $M_\theta$ solo cambia cuando $\theta$ se usa, por lo que el valor de una acción poco
elegida refleja la última vez que se probó y no decae mientras tanto.
\end{itemize}
Las candidatas rechazadas reciben esta misma penalización, igual que las aceptadas que no
mejoran: a diferencia del método clásico, la aceptación sin mejora no se premia.

\subsection{Actualización de la tabla}

El siguiente estado es $z'=1$ si $I_t>0$ y $z'=0$ en caso contrario, y se aplica la regla de
Q-learning~\cite{watkins1992}:
\begin{equation}
Q(z,\theta)\gets Q(z,\theta)+\alpha\Bigl[r_t+\gamma\max_{\theta'}Q(z',\theta')-Q(z,\theta)\Bigr].
\label{eq:qlearning}
\end{equation}
El Algoritmo~\ref{alg:ql-apr} resume el paso de aprendizaje completo, que se ejecuta en
\emph{todas} las iteraciones de las dos fases, incluidas las de calentamiento.

\begin{algorithm}[H]
\caption{Aprendizaje por Q-learning (\cod{ALNS\_QLearning::learn})}
\label{alg:ql-apr}
\Entrada{iteración $t$; resultado $o$; estado $z$ y acción $\theta$ de la iteración}
$I\gets\eta\,\varphi(o.\gbest)+(1-\eta)\,\varphi(o.\gcur)$\;
$\tau\gets t/T_{\mathrm{fin}}$\;
\uSi{$I>0$}{
  $r\gets\tau\,I$;\quad $z'\gets 1$\;
}
\SiNo{
  $\Omega\gets\max_{\theta'\neq\theta}M_{\theta'}$ \Com*[r]{costo de oportunidad}
  $r\gets-\tau\,\Omega$;\quad $z'\gets 0$\;
}
$M_\theta\gets\delta\,M_\theta+(1-\delta)\,I$\;
$Q(z,\theta)\gets Q(z,\theta)+\alpha\bigl[r+\gamma\max_{\theta'}Q(z',\theta')-Q(z,\theta)\bigr]$\;
$z\gets z'$\;
\lSi{$t>t_{\mathrm{exp}}$}{$\varepsilon\gets\max(\varepsilon_{\min},\ \beta\,\varepsilon)$}
\end{algorithm}

\subsection{Reacción al cambio de fase}
\label{sec:ql-fase2}

Al comenzar la fase~2 el objetivo efectivo de la búsqueda cambia: deja de consistir en colocar
clientes pendientes y pasa a ser reducir la distancia de una solución completa. Los pares
que mejor vaciaban el banco no tienen por qué ser los que mejor acortan rutas. El agente
recibe el aviso \fn{AlIniciarFase2} del bucle común (Algoritmo~\ref{alg:etapa}) y responde
como indica el Algoritmo~\ref{alg:ql-fase2}.

\begin{algorithm}[H]
\caption{Inicio de la fase 2 en el Q-learning (\cod{ALNS\_QLearning::onDistancePhase})}
\label{alg:ql-fase2}
\Entrada{iteración $t$ en la que comienza la fase 2}
$M_\theta\gets 0$ para toda acción $\theta$ \Com*[r]{se olvidan las mejoras de la fase 1}
$\varepsilon\gets 1$;\quad $t_{\mathrm{exp}}\gets t+T_{\mathrm{exp}}$ \Com*[r]{se vuelve a explorar}
\end{algorithm}

\begin{itemize}
\item \textbf{Memoria de mejoras a cero.} Las mejoras de la fase~1 están en la escala de los
clientes colocados ($\varphi\approx 7.1$ en el ejemplo anterior, y hasta
$\ln(1+10^{4})\approx 9.2$ al colocar el último cliente pendiente) y no son alcanzables en la
fase~2. Si se conservaran, el costo de oportunidad penalizaría cada iteración sin mejora
con una magnitud que ningún par puede ya compensar.
\item \textbf{Nuevo calendario de exploración.} Se repiten las $200$ iteraciones uniformes y
el decaimiento de $\varepsilon$ de~\eqref{eq:eps}, ahora con $t_0$ igual a la iteración de
inicio de la fase~2, para que todos los pares vuelvan a evaluarse con el nuevo objetivo.
\item \textbf{La tabla $Q$ y el estado se conservan.} Los valores aprendidos en la fase~1 son
el punto de partida de la fase~2; con $\alpha=0.5$ y la nueva exploración uniforme, se
sobrescriben en pocas visitas a cada celda.
\end{itemize}
Las etapas sucesivas de la fase~1, en cambio, comparten objetivo y no provocan ningún
reinicio del agente: la tabla, la memoria de mejoras y $\varepsilon$ continúan de una etapa a
la siguiente.

\subsection{Propiedades relevantes}

\begin{itemize}
\item \textbf{Memoria corta.} Con tasa constante $\alpha=0.5$, el valor $Q(z,\theta)$ es una
media con pesos $2^{-1},2^{-2},\dots$ sobre los objetivos de las últimas visitas a esa celda:
las dos visitas más recientes concentran el 75\,\% del valor. La estimación reacciona de
inmediato a los cambios de régimen de la búsqueda, lo cual es deseable porque el problema es
no estacionario (la utilidad de un operador depende de la fase, de la temperatura y de la
calidad de la solución actual), pero también es ruidosa.
\item \textbf{Recompensa no estacionaria por diseño.} El factor $\tau_t$ hace que la escala de
las recompensas, positivas y negativas, aumente a lo largo de cada fase, y el costo de
oportunidad $\Omega_t$ sube y baja con el rendimiento reciente de los pares; una tasa de
aprendizaje alta es coherente con esa elección.
\item \textbf{Premios y penalizaciones en la misma escala.} La penalización de una iteración
sin mejora es una mejora \emph{media por uso}, no la mayor mejora observada. Para un par que
mejora con frecuencia $p$ y mejora media $\bar{I}$ cuando lo hace, la recompensa esperada
de~\eqref{eq:recompensa} es aproximadamente $\tau\,[\,p\,\bar{I}-(1-p)\,\Omega\,]$, mientras
que su propia memoria tiende a $M_\theta\approx p\,\bar{I}$. En consecuencia, la recompensa
esperada de un par es positiva cuando su mejora media por uso supera, aproximadamente, a la de
la mejor alternativa, y negativa en caso contrario: la señal compara cada par con el mejor de
los demás, combinando frecuencia y magnitud de las mejoras. Esta es una consecuencia analítica
de la fórmula; puede contrastarse empíricamente registrando la evolución de la tabla $Q$, de
$M$ y de la frecuencia de uso de cada par.
\item \textbf{Papel del estado.} El estado distingue la situación «se acaba de mejorar» de
«no se acaba de mejorar». Puesto que las mejoras son minoritarias, $z=0$ es el estado
dominante y su fila es la que se actualiza con más frecuencia; la fila $z=1$ captura qué par
conviene aplicar inmediatamente después de un éxito (por ejemplo, insistir con un operador
intensivo). El término $\gamma\max_{\theta'}Q(z',\theta')$ propaga a cada par el valor de la
situación a la que conduce.
\item \textbf{Concentración.} Con $\varepsilon_{\min}=0.01$, el método concentra casi todas las
iteraciones en el par mejor valorado de cada estado, y cambia de par cuando las
penalizaciones acumuladas lo hacen caer por debajo de otro. El reparto de uso entre operadores
es, por tanto, mucho más desigual y más cambiante que el de la ruleta, que reparte en
proporción a los pesos.
\end{itemize}
